\documentclass[11pt,a4paper]{article}

\usepackage[utf8]{inputenc}
\usepackage[T1]{fontenc}
\usepackage[spanish,es-noshorthands,es-nodecimaldot]{babel}
\usepackage[margin=2.5cm]{geometry}
\usepackage{graphicx}
\usepackage{float}
\usepackage{booktabs}
\usepackage{amsmath}
\usepackage{enumitem}
\usepackage{caption}
\usepackage[hidelinks]{hyperref}

% Las figuras y tablas se generan en ../output/
\graphicspath{{../output/}}

\setlist[itemize]{itemsep=2pt,topsep=4pt}
\captionsetup{font=small,labelfont=bf}

\begin{document}

\begin{titlepage}
\centering
\vspace*{3cm}

{\LARGE\bfseries Trabajo Práctico N$^{\circ}$2\par}
\vspace{0.4cm}
{\large (Semestre de Otoño 2025)\par}
\vspace{0.3cm}
{\large\bfseries Big Data\par}

\vfill

{\large Profesora: María Noelia Romero\par}
\vspace{0.2cm}
{\large Tutor: Tomás Pacheco\par}
\vspace{0.2cm}
{\large Alumnos: Catalina Gallardo, Manuel Elorga\par}

\vfill
\end{titlepage}

% ===============================================================
\section*{Parte I: Limpieza de la base de datos y la maldición de la multidimensionalidad}
% ===============================================================

\subsection*{Ejercicio 1}

La base de datos tiene 168 columnas.

\subsection*{Ejercicio 2}

Si convirtiéramos todas las variables categóricas a \textit{dummies}, la base de datos
tendría 293 columnas.

\subsection*{Ejercicio 3}

Las variables que resultan más relevantes para agrupar a personas desocupadas son:

\begin{itemize}
  \item Tamaño del aglomerado urbano (MAS\_500).
  \item Área urbana (AGLOMERADO).
  \item Sexo (CH04).
  \item Edad (CH06).
  \item Estado civil (CH07).
  \item Cobertura médica (CH08).
  \item Lugar de nacimiento (CH15).
  \item Lugar de residencia hace cinco años (CH16).
  \item Nivel educativo (NIVEL\_ED).
  \item Indica si trabajó en el último año (PP02I).
  \item Indica si trabajó alguna vez (PP10D).
  \item Ingreso per cápita familiar (IPCF).
\end{itemize}

% ===============================================================
\section*{Parte II: Clustering}
% ===============================================================

\subsection*{Ejercicio 2}

\subsubsection*{Inciso a}

Para el análisis de $k$-medias se definieron dos variables centrales:

\begin{itemize}
  \item \textbf{Educación}: se calcularon los años estándar completados según el nivel
        educativo (primaria: 6 años; secundaria: 12 años; terciario: 14 años, etc.),
        ajustando por años aprobados cuando el nivel no estaba finalizado. El resultado es un
        indicador continuo de formación formal.
  \item \textbf{Ingreso}: se actualizaron los valores a pesos constantes del primer trimestre
        de 2024.
\end{itemize}

A partir de estas dos dimensiones se ejecutó $k$-means con tres valores de $K$, como se ve en
la Figura~\ref{fig:kmedias}:

\begin{itemize}
  \item $K = 2$: surge una división clara entre ingresos bajos con menor educación e ingresos
        altos con mayor formación.
  \item $K = 4$: es la segmentación más rica, con grupos que van desde muy bajos ingresos y
        poca educación hasta altos ingresos con formación avanzada.
  \item $K = 10$: se observa una fragmentación fina que captura diferencias sutiles, pero con
        mayor riesgo de ruido y menor interpretabilidad.
\end{itemize}

En resumen, $K = 2$ y $K = 4$ ofrecen agrupaciones claras y útiles, mientras que $K = 10$
aporta más detalle a costa de simplicidad y robustez analítica.

\begin{figure}[H]
  \centering
  \includegraphics[width=\textwidth]{{k-medias}.png}
  \caption{Análisis de $k$-medias con $K=2$, $K=4$ y $K=10$.}
  \label{fig:kmedias}
\end{figure}

\subsubsection*{Inciso b}

El algoritmo $k$-means con $K = 2$ logra separar parcialmente a las personas ocupadas de las
desocupadas, principalmente en función del ingreso y de los años de educación. Se observa que
las personas con mayores ingresos y más educación tienden a agruparse en un mismo clúster,
que generalmente corresponde a los ocupados. Sin embargo, hay un solapamiento considerable:
algunos desocupados tienen características similares a las de los ocupados y viceversa. Esto
muestra que el algoritmo capta ciertas diferencias estructurales, pero no clasifica con
precisión la condición laboral, ya que no considera otros factores relevantes.

\begin{figure}[H]
  \centering
  \includegraphics[width=\textwidth]{{kmedias_comparacion}.png}
  \caption{Análisis de $k$-medias con $K=2$, según condición de ocupación.}
  \label{fig:kmedias2}
\end{figure}

\subsection*{Ejercicio 3}

En este ejercicio se aplica un análisis de clustering jerárquico aglomerativo sobre las
personas desocupadas de la base \texttt{respondieron}, con el objetivo de identificar
subgrupos o patrones homogéneos dentro de la muestra sin necesidad de definir de antemano un
número óptimo de grupos.

Un dendrograma es un diagrama en forma de árbol que muestra cómo, paso a paso, se van uniendo
objetos en un análisis de clustering jerárquico. En el eje vertical se representa la distancia
o disimilitud a la que se fusionan primero pares de observaciones, luego grupos sucesivamente
mayores, hasta formar un único clúster. Así, las uniones más bajas indican alta similitud y
las más altas revelan diferencias mayores entre los grupos.

La Figura~\ref{fig:dendograma} muestra primero fusiones a muy baja distancia, lo que indica
parejas de individuos con perfiles casi idénticos. A distancias intermedias se forman dos
grandes ramas: una reúne sobre todo jóvenes sin experiencia ni cobertura, y otra agrupa
desocupados con algo de trayectoria y seguro médico. Finalmente, dos casos aislados
(\textit{outliers}) se fusionan al resto a muy alta distancia, lo que señala perfiles
atípicos. Un corte horizontal a una distancia intermedia permitiría extraer entre cuatro y
cinco clústeres homogéneos para caracterizar subtipos de desocupación.

\begin{figure}[H]
  \centering
  \includegraphics[width=0.9\textwidth]{{dendograma}.png}
  \caption{Clustering jerárquico.}
  \label{fig:dendograma}
\end{figure}

% ===============================================================
\section*{Parte III: Histogramas y kernels}
% ===============================================================

\subsection*{Ejercicio 1}

La Figura~\ref{fig:histograma} compara la distribución del ingreso individual (P47T) en 2004
y 2024. Para no sesgar el análisis se excluyó de la muestra el cuantil más alto, de modo de
dejar fuera a las personas con ingresos demasiado grandes. No se excluyen los datos del
primer cuantil, ya que la distribución está sesgada hacia la izquierda: hay muchos ingresos
muy bajos.

En ambos histogramas se optó por utilizar 30 bins, para poder identificar todo tipo de picos,
asimetrías y modas de los ingresos. Dado que tenemos miles de observaciones, esa cantidad
asegura que cada bin contenga un número significativo de casos. Además, se utilizó la misma
cantidad de bins para ambos años, para facilitar la comparación.

Se observa que la distribución de 2004 presenta una mayor cantidad de ingresos bajos. En
cambio, en 2024 la distribución se orienta hacia la derecha, lo que refleja un incremento de
los ingresos individuales. Además, se puede observar una mayor heterogeneidad de los datos, ya
que hay una gran dispersión con una mayor cantidad de valores.

\begin{figure}[H]
  \centering
  \includegraphics[width=\textwidth]{{histograma}.png}
  \caption{Distribución del ingreso individual, 2004 y 2024.}
  \label{fig:histograma}
\end{figure}

\subsection*{Ejercicio 2}

Para este ejercicio se creó una nueva base de datos llamada \texttt{hogares}, utilizando la
variable CODUSU que identifica a cada hogar. En esa base se agregaron la variable binaria
\texttt{hay\_desocupado}, que indica si hay al menos un desocupado en el hogar, y la variable
\texttt{ingreso\_hogar}, que indica el ingreso total del hogar.

\subsubsection*{Inciso c}

En la Figura~\ref{fig:unif2004} (2004), los hogares con al menos una persona desocupada
(línea roja) concentran su densidad en ingresos bajos, con mayor variabilidad reflejada en
picos pequeños. En contraste, los hogares con todos sus miembros ocupados (línea verde)
presentan ingresos promedio más altos, con una distribución más uniforme y menos
oscilaciones.

En la Figura~\ref{fig:unif2024} (2024), los hogares con desocupados mantienen su
concentración en ingresos bajos, pero con un rango más amplio y mayor dispersión, lo que
indica una desigualdad interna persistente. Los hogares con todos ocupados siguen mostrando
ingresos más altos y un incremento general respecto de 2004, aunque la diferencia con los
hogares con desocupados se mantiene marcada.

\begin{figure}[H]
  \centering
  \begin{minipage}{0.49\textwidth}
    \centering
    \includegraphics[width=\textwidth]{{distribucion2004}.png}
    \caption{Ingreso por hogar con kernel uniforme, 2004.}
    \label{fig:unif2004}
  \end{minipage}\hfill
  \begin{minipage}{0.49\textwidth}
    \centering
    \includegraphics[width=\textwidth]{{distribucion2024}.png}
    \caption{Ingreso por hogar con kernel uniforme, 2024.}
    \label{fig:unif2024}
  \end{minipage}
\end{figure}

\subsubsection*{Inciso d}

El ancho de banda es el tamaño de la ``ventana'' que utilizamos para calcular cuánto influyen
unos datos sobre otros, de manera que el resultado sea lo más preciso posible. Para elegir
cuál es el mejor ancho de banda debemos considerar un \textit{trade-off} entre sesgo y
varianza.

Si elegimos un ancho de banda muy grande, estamos promediando demasiados puntos y se pierden
detalles importantes. En cambio, si el ancho de banda es muy pequeño, la estimación se ajusta
demasiado a cada punto, lo que significa que se concentra en exceso en las variaciones
individuales.

El ancho de banda óptimo es ese valor intermedio en el que compensamos la pérdida de detalle
(sesgo) y el exceso de variaciones (varianza), para aproximarnos lo mejor posible a la
verdadera forma de los datos.

\subsubsection*{Inciso e}

La diferencia porcentual en ingresos entre hogares con todos sus miembros ocupados y hogares
con desocupados disminuyó del 38\% en 2004 al 27\% en 2024, lo que indica una reducción en la
brecha relativa entre ambos grupos. Esto podría reflejar mejoras en los ingresos de los
hogares con desocupados, un crecimiento más lento en los ingresos de los hogares ocupados, o
cambios en las dinámicas laborales y sociales que mitigaron parcialmente el impacto de la
desocupación. Sin embargo, aunque la brecha relativa se redujo, sigue siendo significativa.

\subsection*{Ejercicio 3}

\subsubsection*{Inciso a}

En la Figura~\ref{fig:gauss2004}, la curva de los hogares con desocupados (línea roja) se
encuentra desplazada hacia la izquierda respecto de la de los hogares sin desocupados (línea
verde). Esto sugiere que los primeros presentan ingresos más bajos. A su vez, los hogares sin
desocupados muestran una mayor dispersión hacia ingresos más elevados.

En la Figura~\ref{fig:gauss2024}, los hogares con desocupados tienden a concentrarse en los
niveles más bajos de ingreso. Sin embargo, podemos ver que ambas distribuciones aparecen más
suavizadas y desplazadas hacia la derecha respecto de 2004, lo que refleja un aumento general
de los ingresos.

\begin{figure}[H]
  \centering
  \begin{minipage}{0.49\textwidth}
    \centering
    \includegraphics[width=\textwidth]{{distribuciong2004}.png}
    \caption{Ingreso por hogar con kernel gaussiano, 2004.}
    \label{fig:gauss2004}
  \end{minipage}\hfill
  \begin{minipage}{0.49\textwidth}
    \centering
    \includegraphics[width=\textwidth]{{distribuciong2024}.png}
    \caption{Ingreso por hogar con kernel gaussiano, 2024.}
    \label{fig:gauss2024}
  \end{minipage}
\end{figure}

En la Figura~\ref{fig:epa2004}, el uso del kernel de Epanechnikov genera picos más marcados y
cambios más bruscos en la densidad. Igualmente, llegamos a la misma conclusión: los hogares
con desocupados tienen menores ingresos, con una curva más centrada a la izquierda.

En la Figura~\ref{fig:epa2024}, la estimación es más irregular que la gaussiana. Nuevamente,
los hogares sin desocupados presentan ingresos más altos y una densidad más dispersa,
mientras que los hogares con desocupados se concentran en los rangos inferiores de ingresos.

\begin{figure}[H]
  \centering
  \begin{minipage}{0.49\textwidth}
    \centering
    \includegraphics[width=\textwidth]{{distribucione2004}.png}
    \caption{Ingreso por hogar con kernel de Epanechnikov, 2004.}
    \label{fig:epa2004}
  \end{minipage}\hfill
  \begin{minipage}{0.49\textwidth}
    \centering
    \includegraphics[width=\textwidth]{{distribucione2024}.png}
    \caption{Ingreso por hogar con kernel de Epanechnikov, 2024.}
    \label{fig:epa2024}
  \end{minipage}
\end{figure}

\subsubsection*{Inciso b}

El kernel gaussiano genera curvas más suaves y continuas, ya que asigna peso a todos los
puntos del dominio. Permite capturar la tendencia general de la distribución del ingreso sin
cambios abruptos y muestra colas más prolongadas.

Por otro lado, el kernel de Epanechnikov produce colas más segmentadas y con cambios más
abruptos. Es por esto que la estimación resulta más sensible a pequeñas concentraciones de
datos, en las que se visualizan picos más marcados.

\end{document}
